\section{Introduction}
Reinforcement learning (RL) formulates control as sequential decision making, where an agent learns a policy by interacting with an environment and maximizing expected return \cite{sutton2018reinforcement}. In robotics this setting is attractive because it can optimize behaviors that are difficult to hand-design, but it is also difficult because interactions are expensive, exploration can be unsafe, and the learned policy can exploit simulator-specific details \cite{kober2013reinforcement,kormushev2013reinforcement}. A common practical solution is to train in simulation and transfer the policy to the real robot. The central difficulty is the reality gap: the simulator dynamics are never exactly equal to the deployment dynamics \cite{hofer2020perspectives}.

This project follows the course sim-to-real abstraction in a controlled sim-to-sim benchmark. In Part 1, we implement basic policy-gradient algorithms on \texttt{Hopper-v4} to study the variance and stability issues behind direct policy optimization. In Part 2, we consider \texttt{PandaPush-v3}, where the source domain involves a light object and the target domain involves a heavier object. The target domain serves as a proxy for the real world: it is available for evaluation, but direct training on it is treated as an upper bound, rather than as a practical deployable solution.

The second part uses modern actor-critic algorithms from Stable-Baselines3. In particular, we consider PPO, which constrains policy updates through a clipped surrogate objective \cite{schulman2017ppo}, and SAC, an off-policy maximum-entropy method that learns stochastic policies and reuses data through a replay buffer \cite{haarnoja2018sac}. To reduce the transfer gap, we also implement domain randomization, which trains the agent over a distribution of simulator parameters rather than a single nominal model \cite{tobin2017domain,peng2018dynamics,muratore2022robot}. Specifically, the implemented UDR wrapper samples parameters from fixed ranges, while the ADR wrapper adapts the mass range based on recent success, following the curriculum idea behind automatic domain randomization \cite{akkaya2019solving}. These algorithms and randomization strategies will be discussed in more detail in the following chapters.

